% Options for packages loaded elsewhere
\PassOptionsToPackage{unicode}{hyperref}
\PassOptionsToPackage{hyphens}{url}
\documentclass[
]{article}
\usepackage{xcolor}
\usepackage[margin=1in]{geometry}
\usepackage{amsmath,amssymb}
\setcounter{secnumdepth}{-\maxdimen} % remove section numbering
\usepackage{iftex}
\ifPDFTeX
  \usepackage[T1]{fontenc}
  \usepackage[utf8]{inputenc}
  \usepackage{textcomp} % provide euro and other symbols
\else % if luatex or xetex
  \usepackage{unicode-math} % this also loads fontspec
  \defaultfontfeatures{Scale=MatchLowercase}
  \defaultfontfeatures[\rmfamily]{Ligatures=TeX,Scale=1}
\fi
\usepackage{lmodern}
\ifPDFTeX\else
  % xetex/luatex font selection
\fi
% Use upquote if available, for straight quotes in verbatim environments
\IfFileExists{upquote.sty}{\usepackage{upquote}}{}
\IfFileExists{microtype.sty}{% use microtype if available
  \usepackage[]{microtype}
  \UseMicrotypeSet[protrusion]{basicmath} % disable protrusion for tt fonts
}{}
\makeatletter
\@ifundefined{KOMAClassName}{% if non-KOMA class
  \IfFileExists{parskip.sty}{%
    \usepackage{parskip}
  }{% else
    \setlength{\parindent}{0pt}
    \setlength{\parskip}{6pt plus 2pt minus 1pt}}
}{% if KOMA class
  \KOMAoptions{parskip=half}}
\makeatother
\usepackage{color}
\usepackage{fancyvrb}
\newcommand{\VerbBar}{|}
\newcommand{\VERB}{\Verb[commandchars=\\\{\}]}
\DefineVerbatimEnvironment{Highlighting}{Verbatim}{commandchars=\\\{\}}
% Add ',fontsize=\small' for more characters per line
\usepackage{framed}
\definecolor{shadecolor}{RGB}{248,248,248}
\newenvironment{Shaded}{\begin{snugshade}}{\end{snugshade}}
\newcommand{\AlertTok}[1]{\textcolor[rgb]{0.94,0.16,0.16}{#1}}
\newcommand{\AnnotationTok}[1]{\textcolor[rgb]{0.56,0.35,0.01}{\textbf{\textit{#1}}}}
\newcommand{\AttributeTok}[1]{\textcolor[rgb]{0.13,0.29,0.53}{#1}}
\newcommand{\BaseNTok}[1]{\textcolor[rgb]{0.00,0.00,0.81}{#1}}
\newcommand{\BuiltInTok}[1]{#1}
\newcommand{\CharTok}[1]{\textcolor[rgb]{0.31,0.60,0.02}{#1}}
\newcommand{\CommentTok}[1]{\textcolor[rgb]{0.56,0.35,0.01}{\textit{#1}}}
\newcommand{\CommentVarTok}[1]{\textcolor[rgb]{0.56,0.35,0.01}{\textbf{\textit{#1}}}}
\newcommand{\ConstantTok}[1]{\textcolor[rgb]{0.56,0.35,0.01}{#1}}
\newcommand{\ControlFlowTok}[1]{\textcolor[rgb]{0.13,0.29,0.53}{\textbf{#1}}}
\newcommand{\DataTypeTok}[1]{\textcolor[rgb]{0.13,0.29,0.53}{#1}}
\newcommand{\DecValTok}[1]{\textcolor[rgb]{0.00,0.00,0.81}{#1}}
\newcommand{\DocumentationTok}[1]{\textcolor[rgb]{0.56,0.35,0.01}{\textbf{\textit{#1}}}}
\newcommand{\ErrorTok}[1]{\textcolor[rgb]{0.64,0.00,0.00}{\textbf{#1}}}
\newcommand{\ExtensionTok}[1]{#1}
\newcommand{\FloatTok}[1]{\textcolor[rgb]{0.00,0.00,0.81}{#1}}
\newcommand{\FunctionTok}[1]{\textcolor[rgb]{0.13,0.29,0.53}{\textbf{#1}}}
\newcommand{\ImportTok}[1]{#1}
\newcommand{\InformationTok}[1]{\textcolor[rgb]{0.56,0.35,0.01}{\textbf{\textit{#1}}}}
\newcommand{\KeywordTok}[1]{\textcolor[rgb]{0.13,0.29,0.53}{\textbf{#1}}}
\newcommand{\NormalTok}[1]{#1}
\newcommand{\OperatorTok}[1]{\textcolor[rgb]{0.81,0.36,0.00}{\textbf{#1}}}
\newcommand{\OtherTok}[1]{\textcolor[rgb]{0.56,0.35,0.01}{#1}}
\newcommand{\PreprocessorTok}[1]{\textcolor[rgb]{0.56,0.35,0.01}{\textit{#1}}}
\newcommand{\RegionMarkerTok}[1]{#1}
\newcommand{\SpecialCharTok}[1]{\textcolor[rgb]{0.81,0.36,0.00}{\textbf{#1}}}
\newcommand{\SpecialStringTok}[1]{\textcolor[rgb]{0.31,0.60,0.02}{#1}}
\newcommand{\StringTok}[1]{\textcolor[rgb]{0.31,0.60,0.02}{#1}}
\newcommand{\VariableTok}[1]{\textcolor[rgb]{0.00,0.00,0.00}{#1}}
\newcommand{\VerbatimStringTok}[1]{\textcolor[rgb]{0.31,0.60,0.02}{#1}}
\newcommand{\WarningTok}[1]{\textcolor[rgb]{0.56,0.35,0.01}{\textbf{\textit{#1}}}}
\usepackage{graphicx}
\makeatletter
\newsavebox\pandoc@box
\newcommand*\pandocbounded[1]{% scales image to fit in text height/width
  \sbox\pandoc@box{#1}%
  \Gscale@div\@tempa{\textheight}{\dimexpr\ht\pandoc@box+\dp\pandoc@box\relax}%
  \Gscale@div\@tempb{\linewidth}{\wd\pandoc@box}%
  \ifdim\@tempb\p@<\@tempa\p@\let\@tempa\@tempb\fi% select the smaller of both
  \ifdim\@tempa\p@<\p@\scalebox{\@tempa}{\usebox\pandoc@box}%
  \else\usebox{\pandoc@box}%
  \fi%
}
% Set default figure placement to htbp
\def\fps@figure{htbp}
\makeatother
\setlength{\emergencystretch}{3em} % prevent overfull lines
\providecommand{\tightlist}{%
  \setlength{\itemsep}{0pt}\setlength{\parskip}{0pt}}
\usepackage{bookmark}
\IfFileExists{xurl.sty}{\usepackage{xurl}}{} % add URL line breaks if available
\urlstyle{same}
\hypersetup{
  pdftitle={750},
  hidelinks,
  pdfcreator={LaTeX via pandoc}}

\title{750}
\author{}
\date{\vspace{-2.5em}2026-05-20}

\begin{document}
\maketitle

\subsection{R Markdown}\label{r-markdown}

This is an R Markdown document. Markdown is a simple formatting syntax
for authoring HTML, PDF, and MS Word documents. For more details on
using R Markdown see \url{http://rmarkdown.rstudio.com}.

When you click the \textbf{Knit} button a document will be generated
that includes both content as well as the output of any embedded R code
chunks within the document. You can embed an R code chunk like this:

\begin{Shaded}
\begin{Highlighting}[]
\FunctionTok{library}\NormalTok{(DESeq2)}
\end{Highlighting}
\end{Shaded}

\begin{verbatim}
## Loading required package: S4Vectors
\end{verbatim}

\begin{verbatim}
## Loading required package: stats4
\end{verbatim}

\begin{verbatim}
## Loading required package: BiocGenerics
\end{verbatim}

\begin{verbatim}
## Loading required package: generics
\end{verbatim}

\begin{verbatim}
## 
## Attaching package: 'generics'
\end{verbatim}

\begin{verbatim}
## The following objects are masked from 'package:base':
## 
##     as.difftime, as.factor, as.ordered, intersect, is.element, setdiff,
##     setequal, union
\end{verbatim}

\begin{verbatim}
## 
## Attaching package: 'BiocGenerics'
\end{verbatim}

\begin{verbatim}
## The following objects are masked from 'package:stats':
## 
##     IQR, mad, sd, var, xtabs
\end{verbatim}

\begin{verbatim}
## The following objects are masked from 'package:base':
## 
##     anyDuplicated, aperm, append, as.data.frame, basename, cbind,
##     colnames, dirname, do.call, duplicated, eval, evalq, Filter, Find,
##     get, grep, grepl, is.unsorted, lapply, Map, mapply, match, mget,
##     order, paste, pmax, pmax.int, pmin, pmin.int, Position, rank,
##     rbind, Reduce, rownames, sapply, saveRDS, table, tapply, unique,
##     unsplit, which.max, which.min
\end{verbatim}

\begin{verbatim}
## 
## Attaching package: 'S4Vectors'
\end{verbatim}

\begin{verbatim}
## The following object is masked from 'package:utils':
## 
##     findMatches
\end{verbatim}

\begin{verbatim}
## The following objects are masked from 'package:base':
## 
##     expand.grid, I, unname
\end{verbatim}

\begin{verbatim}
## Loading required package: IRanges
\end{verbatim}

\begin{verbatim}
## 
## Attaching package: 'IRanges'
\end{verbatim}

\begin{verbatim}
## The following object is masked from 'package:grDevices':
## 
##     windows
\end{verbatim}

\begin{verbatim}
## Loading required package: GenomicRanges
\end{verbatim}

\begin{verbatim}
## Loading required package: Seqinfo
\end{verbatim}

\begin{verbatim}
## Loading required package: SummarizedExperiment
\end{verbatim}

\begin{verbatim}
## Loading required package: MatrixGenerics
\end{verbatim}

\begin{verbatim}
## Loading required package: matrixStats
\end{verbatim}

\begin{verbatim}
## 
## Attaching package: 'MatrixGenerics'
\end{verbatim}

\begin{verbatim}
## The following objects are masked from 'package:matrixStats':
## 
##     colAlls, colAnyNAs, colAnys, colAvgsPerRowSet, colCollapse,
##     colCounts, colCummaxs, colCummins, colCumprods, colCumsums,
##     colDiffs, colIQRDiffs, colIQRs, colLogSumExps, colMadDiffs,
##     colMads, colMaxs, colMeans2, colMedians, colMins, colOrderStats,
##     colProds, colQuantiles, colRanges, colRanks, colSdDiffs, colSds,
##     colSums2, colTabulates, colVarDiffs, colVars, colWeightedMads,
##     colWeightedMeans, colWeightedMedians, colWeightedSds,
##     colWeightedVars, rowAlls, rowAnyNAs, rowAnys, rowAvgsPerColSet,
##     rowCollapse, rowCounts, rowCummaxs, rowCummins, rowCumprods,
##     rowCumsums, rowDiffs, rowIQRDiffs, rowIQRs, rowLogSumExps,
##     rowMadDiffs, rowMads, rowMaxs, rowMeans2, rowMedians, rowMins,
##     rowOrderStats, rowProds, rowQuantiles, rowRanges, rowRanks,
##     rowSdDiffs, rowSds, rowSums2, rowTabulates, rowVarDiffs, rowVars,
##     rowWeightedMads, rowWeightedMeans, rowWeightedMedians,
##     rowWeightedSds, rowWeightedVars
\end{verbatim}

\begin{verbatim}
## Loading required package: Biobase
\end{verbatim}

\begin{verbatim}
## Warning: package 'Biobase' was built under R version 4.5.3
\end{verbatim}

\begin{verbatim}
## Welcome to Bioconductor
## 
##     Vignettes contain introductory material; view with
##     'browseVignettes()'. To cite Bioconductor, see
##     'citation("Biobase")', and for packages 'citation("pkgname")'.
\end{verbatim}

\begin{verbatim}
## 
## Attaching package: 'Biobase'
\end{verbatim}

\begin{verbatim}
## The following object is masked from 'package:MatrixGenerics':
## 
##     rowMedians
\end{verbatim}

\begin{verbatim}
## The following objects are masked from 'package:matrixStats':
## 
##     anyMissing, rowMedians
\end{verbatim}

\begin{Shaded}
\begin{Highlighting}[]
\FunctionTok{library}\NormalTok{(tidyverse)}
\end{Highlighting}
\end{Shaded}

\begin{verbatim}
## Warning: package 'ggplot2' was built under R version 4.5.3
\end{verbatim}

\begin{verbatim}
## Warning: package 'readr' was built under R version 4.5.3
\end{verbatim}

\begin{verbatim}
## -- Attaching core tidyverse packages ------------------------ tidyverse 2.0.0 --
## v dplyr     1.1.4     v readr     2.2.0
## v forcats   1.0.1     v stringr   1.6.0
## v ggplot2   4.0.3     v tibble    3.3.1
## v lubridate 1.9.4     v tidyr     1.3.2
## v purrr     1.2.1
\end{verbatim}

\begin{verbatim}
## -- Conflicts ------------------------------------------ tidyverse_conflicts() --
## x lubridate::%within%() masks IRanges::%within%()
## x dplyr::collapse()     masks IRanges::collapse()
## x dplyr::combine()      masks Biobase::combine(), BiocGenerics::combine()
## x dplyr::count()        masks matrixStats::count()
## x dplyr::desc()         masks IRanges::desc()
## x tidyr::expand()       masks S4Vectors::expand()
## x dplyr::filter()       masks stats::filter()
## x dplyr::first()        masks S4Vectors::first()
## x dplyr::lag()          masks stats::lag()
## x ggplot2::Position()   masks BiocGenerics::Position(), base::Position()
## x purrr::reduce()       masks GenomicRanges::reduce(), IRanges::reduce()
## x dplyr::rename()       masks S4Vectors::rename()
## x lubridate::second()   masks S4Vectors::second()
## x lubridate::second<-() masks S4Vectors::second<-()
## x dplyr::slice()        masks IRanges::slice()
## i Use the conflicted package (<http://conflicted.r-lib.org/>) to force all conflicts to become errors
\end{verbatim}

\begin{Shaded}
\begin{Highlighting}[]
\FunctionTok{citation}\NormalTok{(}\StringTok{"tidyverse"}\NormalTok{)}
\end{Highlighting}
\end{Shaded}

\begin{verbatim}
## To cite package 'tidyverse' in publications use:
## 
##   Wickham H, Averick M, Bryan J, Chang W, McGowan LD, François R,
##   Grolemund G, Hayes A, Henry L, Hester J, Kuhn M, Pedersen TL, Miller
##   E, Bache SM, Müller K, Ooms J, Robinson D, Seidel DP, Spinu V,
##   Takahashi K, Vaughan D, Wilke C, Woo K, Yutani H (2019). "Welcome to
##   the tidyverse." _Journal of Open Source Software_, *4*(43), 1686.
##   doi:10.21105/joss.01686 <https://doi.org/10.21105/joss.01686>.
## 
## A BibTeX entry for LaTeX users is
## 
##   @Article{,
##     title = {Welcome to the {tidyverse}},
##     author = {Hadley Wickham and Mara Averick and Jennifer Bryan and Winston Chang and Lucy D'Agostino McGowan and Romain François and Garrett Grolemund and Alex Hayes and Lionel Henry and Jim Hester and Max Kuhn and Thomas Lin Pedersen and Evan Miller and Stephan Milton Bache and Kirill Müller and Jeroen Ooms and David Robinson and Dana Paige Seidel and Vitalie Spinu and Kohske Takahashi and Davis Vaughan and Claus Wilke and Kara Woo and Hiroaki Yutani},
##     year = {2019},
##     journal = {Journal of Open Source Software},
##     volume = {4},
##     number = {43},
##     pages = {1686},
##     doi = {10.21105/joss.01686},
##   }
\end{verbatim}

\begin{Shaded}
\begin{Highlighting}[]
\NormalTok{counts\_data }\OtherTok{\textless{}{-}} \FunctionTok{read.csv}\NormalTok{(}\StringTok{"gene\_counts.tsv"}\NormalTok{,}
                        \AttributeTok{sep=}\StringTok{"}\SpecialCharTok{\textbackslash{}t}\StringTok{"}\NormalTok{, }\AttributeTok{row.names=}\DecValTok{1}\NormalTok{)}
\FunctionTok{head}\NormalTok{(counts\_data)}
\end{Highlighting}
\end{Shaded}

\begin{verbatim}
##           mutant_1 mutant_2 mutant_3 mutant_4 mutant_5 normal_1 normal_2
## GENE_0001       70       81       98       46      108      119      100
## GENE_0002      149      231      298      103      199      206      264
## GENE_0003       31       46       25       29       39       24       21
## GENE_0004      413      438      634      670      617      610      602
## GENE_0005       61       69       77       42       74       91       70
## GENE_0006       26       18       32        5       17       32       25
##           normal_3 normal_4 normal_5
## GENE_0001       66       60       73
## GENE_0002      194      148      196
## GENE_0003       19       26       46
## GENE_0004      408      490      907
## GENE_0005       62       87      112
## GENE_0006       22       17       16
\end{verbatim}

\begin{Shaded}
\begin{Highlighting}[]
\CommentTok{\# Create sample info}
\NormalTok{colData }\OtherTok{\textless{}{-}} \FunctionTok{data.frame}\NormalTok{(}
  \AttributeTok{condition =} \FunctionTok{factor}\NormalTok{(}\FunctionTok{c}\NormalTok{(}\FunctionTok{rep}\NormalTok{(}\StringTok{"normal"}\NormalTok{,}\DecValTok{5}\NormalTok{), }\FunctionTok{rep}\NormalTok{(}\StringTok{"mutant"}\NormalTok{,}\DecValTok{5}\NormalTok{)))}
\NormalTok{)}
\FunctionTok{rownames}\NormalTok{(colData) }\OtherTok{\textless{}{-}} \FunctionTok{colnames}\NormalTok{(counts\_data)}

\CommentTok{\# Check they match {-} must say TRUE}
\FunctionTok{all}\NormalTok{(}\FunctionTok{colnames}\NormalTok{(counts\_data) }\SpecialCharTok{\%in\%} \FunctionTok{rownames}\NormalTok{(colData))}
\end{Highlighting}
\end{Shaded}

\begin{verbatim}
## [1] TRUE
\end{verbatim}

\begin{Shaded}
\begin{Highlighting}[]
\FunctionTok{all}\NormalTok{(}\FunctionTok{colnames}\NormalTok{(counts\_data) }\SpecialCharTok{==} \FunctionTok{rownames}\NormalTok{(colData))}
\end{Highlighting}
\end{Shaded}

\begin{verbatim}
## [1] TRUE
\end{verbatim}

\begin{Shaded}
\begin{Highlighting}[]
\CommentTok{\# Build DESeq2 object}
\NormalTok{dds }\OtherTok{\textless{}{-}} \FunctionTok{DESeqDataSetFromMatrix}\NormalTok{(}\AttributeTok{countData =}\NormalTok{ counts\_data,}
                              \AttributeTok{colData =}\NormalTok{ colData,}
                              \AttributeTok{design =} \SpecialCharTok{\textasciitilde{}}\NormalTok{ condition)}

\CommentTok{\# Filter low counts}
\NormalTok{keep }\OtherTok{\textless{}{-}} \FunctionTok{rowSums}\NormalTok{(}\FunctionTok{counts}\NormalTok{(dds)) }\SpecialCharTok{\textgreater{}=} \DecValTok{10}
\NormalTok{dds }\OtherTok{\textless{}{-}}\NormalTok{ dds[keep,]}

\CommentTok{\# Set normal as reference}
\NormalTok{dds}\SpecialCharTok{$}\NormalTok{condition }\OtherTok{\textless{}{-}} \FunctionTok{relevel}\NormalTok{(dds}\SpecialCharTok{$}\NormalTok{condition, }\AttributeTok{ref =} \StringTok{"normal"}\NormalTok{)}

\CommentTok{\# Run DESeq2}
\NormalTok{dds }\OtherTok{\textless{}{-}} \FunctionTok{DESeq}\NormalTok{(dds)}
\end{Highlighting}
\end{Shaded}

\begin{verbatim}
## estimating size factors
\end{verbatim}

\begin{verbatim}
## estimating dispersions
\end{verbatim}

\begin{verbatim}
## gene-wise dispersion estimates
\end{verbatim}

\begin{verbatim}
## mean-dispersion relationship
\end{verbatim}

\begin{verbatim}
## final dispersion estimates
\end{verbatim}

\begin{verbatim}
## fitting model and testing
\end{verbatim}

\begin{Shaded}
\begin{Highlighting}[]
\NormalTok{res }\OtherTok{\textless{}{-}} \FunctionTok{results}\NormalTok{(dds)}
\NormalTok{res }\OtherTok{\textless{}{-}}\NormalTok{ res[}\FunctionTok{order}\NormalTok{(res}\SpecialCharTok{$}\NormalTok{padj),]}
\FunctionTok{summary}\NormalTok{(res)}
\end{Highlighting}
\end{Shaded}

\begin{verbatim}
## 
## out of 1000 with nonzero total read count
## adjusted p-value < 0.1
## LFC > 0 (up)       : 44, 4.4%
## LFC < 0 (down)     : 51, 5.1%
## outliers [1]       : 0, 0%
## low counts [2]     : 0, 0%
## (mean count < 17)
## [1] see 'cooksCutoff' argument of ?results
## [2] see 'independentFiltering' argument of ?results
\end{verbatim}

\begin{Shaded}
\begin{Highlighting}[]
\CommentTok{\# Save all results}
\NormalTok{res\_df }\OtherTok{\textless{}{-}} \FunctionTok{as.data.frame}\NormalTok{(res)}
\FunctionTok{write.csv}\NormalTok{(res\_df,}
  \StringTok{"DESeq2\_results.csv"}\NormalTok{)}
\CommentTok{\# Save significant DEGs only}
\NormalTok{sig }\OtherTok{\textless{}{-}} \FunctionTok{subset}\NormalTok{(res\_df, padj }\SpecialCharTok{\textless{}} \FloatTok{0.05} \SpecialCharTok{\&} \FunctionTok{abs}\NormalTok{(log2FoldChange) }\SpecialCharTok{\textgreater{}} \DecValTok{1}\NormalTok{)}
\FunctionTok{write.csv}\NormalTok{(sig,}
  \StringTok{"significant\_DEGs.csv"}\NormalTok{)}
\end{Highlighting}
\end{Shaded}

\begin{Shaded}
\begin{Highlighting}[]
\NormalTok{up\_count }\OtherTok{\textless{}{-}} \FunctionTok{sum}\NormalTok{(res\_df}\SpecialCharTok{$}\NormalTok{padj }\SpecialCharTok{\textless{}} \FloatTok{0.05} \SpecialCharTok{\&}\NormalTok{ res\_df}\SpecialCharTok{$}\NormalTok{log2FoldChange }\SpecialCharTok{\textgreater{}} \DecValTok{1}\NormalTok{, }\AttributeTok{na.rm =} \ConstantTok{TRUE}\NormalTok{)}
\FunctionTok{print}\NormalTok{(}\FunctionTok{paste}\NormalTok{(}\StringTok{"Upregulated genes:"}\NormalTok{, up\_count)) }
\end{Highlighting}
\end{Shaded}

\begin{verbatim}
## [1] "Upregulated genes: 19"
\end{verbatim}

\begin{Shaded}
\begin{Highlighting}[]
\NormalTok{down\_count }\OtherTok{\textless{}{-}} \FunctionTok{sum}\NormalTok{(res\_df}\SpecialCharTok{$}\NormalTok{padj }\SpecialCharTok{\textless{}} \FloatTok{0.05} \SpecialCharTok{\&}\NormalTok{ res\_df}\SpecialCharTok{$}\NormalTok{log2FoldChange }\SpecialCharTok{\textless{}} \SpecialCharTok{{-}}\DecValTok{1}\NormalTok{, }\AttributeTok{na.rm =} \ConstantTok{TRUE}\NormalTok{)}
\FunctionTok{print}\NormalTok{(}\FunctionTok{paste}\NormalTok{(}\StringTok{"Downregulated genes:"}\NormalTok{, down\_count))}
\end{Highlighting}
\end{Shaded}

\begin{verbatim}
## [1] "Downregulated genes: 27"
\end{verbatim}

\begin{Shaded}
\begin{Highlighting}[]
\NormalTok{upregulated\_genes }\OtherTok{\textless{}{-}} \FunctionTok{rownames}\NormalTok{(res\_df[}\FunctionTok{which}\NormalTok{(res\_df}\SpecialCharTok{$}\NormalTok{padj }\SpecialCharTok{\textless{}} \FloatTok{0.05} \SpecialCharTok{\&}\NormalTok{ res\_df}\SpecialCharTok{$}\NormalTok{log2FoldChange }\SpecialCharTok{\textgreater{}} \DecValTok{1}\NormalTok{), ])}

\NormalTok{downregulated\_genes }\OtherTok{\textless{}{-}} \FunctionTok{rownames}\NormalTok{(res\_df[}\FunctionTok{which}\NormalTok{(res\_df}\SpecialCharTok{$}\NormalTok{padj }\SpecialCharTok{\textless{}} \FloatTok{0.05} \SpecialCharTok{\&}\NormalTok{ res\_df}\SpecialCharTok{$}\NormalTok{log2FoldChange }\SpecialCharTok{\textless{}} \SpecialCharTok{{-}}\DecValTok{1}\NormalTok{), ])}

\FunctionTok{print}\NormalTok{(upregulated\_genes)}
\end{Highlighting}
\end{Shaded}

\begin{verbatim}
##  [1] "GENE_0098" "GENE_0851" "GENE_0097" "GENE_0496" "GENE_0792" "GENE_0760"
##  [7] "GENE_0227" "GENE_0409" "GENE_0869" "GENE_0237" "GENE_0034" "GENE_0705"
## [13] "GENE_0582" "GENE_0715" "GENE_0658" "GENE_0556" "GENE_0039" "GENE_0448"
## [19] "GENE_0677"
\end{verbatim}

\begin{Shaded}
\begin{Highlighting}[]
\FunctionTok{print}\NormalTok{(downregulated\_genes)}
\end{Highlighting}
\end{Shaded}

\begin{verbatim}
##  [1] "GENE_0818" "GENE_0234" "GENE_0107" "GENE_0199" "GENE_0336" "GENE_0638"
##  [7] "GENE_0680" "GENE_0659" "GENE_0366" "GENE_0457" "GENE_0177" "GENE_0786"
## [13] "GENE_0075" "GENE_0162" "GENE_0239" "GENE_0321" "GENE_0058" "GENE_0780"
## [19] "GENE_0087" "GENE_0408" "GENE_0504" "GENE_0140" "GENE_0611" "GENE_0791"
## [25] "GENE_0225" "GENE_0053" "GENE_0188"
\end{verbatim}

\begin{Shaded}
\begin{Highlighting}[]
\CommentTok{\# Plots (save these for your report)}
\NormalTok{vsd }\OtherTok{\textless{}{-}} \FunctionTok{vst}\NormalTok{(dds, }\AttributeTok{blind=}\ConstantTok{FALSE}\NormalTok{)}
\FunctionTok{plotPCA}\NormalTok{(vsd, }\AttributeTok{intgroup=}\StringTok{"condition"}\NormalTok{)   }
\end{Highlighting}
\end{Shaded}

\begin{verbatim}
## using ntop=500 top features by variance
\end{verbatim}

\pandocbounded{\includegraphics[keepaspectratio]{750_files/figure-latex/unnamed-chunk-9-1.pdf}}

\begin{Shaded}
\begin{Highlighting}[]
\FunctionTok{plotMA}\NormalTok{(res, }\AttributeTok{ylim=}\FunctionTok{c}\NormalTok{(}\SpecialCharTok{{-}}\DecValTok{5}\NormalTok{,}\DecValTok{5}\NormalTok{))            }
\end{Highlighting}
\end{Shaded}

\pandocbounded{\includegraphics[keepaspectratio]{750_files/figure-latex/unnamed-chunk-9-2.pdf}}

\begin{Shaded}
\begin{Highlighting}[]
\CommentTok{\# Load bound genes}
\NormalTok{bound\_genes }\OtherTok{\textless{}{-}} \FunctionTok{read.table}\NormalTok{(}
  \StringTok{"bound\_gene\_ids.txt"}\NormalTok{,}
  \AttributeTok{header=}\ConstantTok{FALSE}\NormalTok{, }\AttributeTok{stringsAsFactors=}\ConstantTok{FALSE}\NormalTok{)}
\NormalTok{bound\_genes}\SpecialCharTok{$}\NormalTok{V1 }\OtherTok{\textless{}{-}} \FunctionTok{gsub}\NormalTok{(}\StringTok{"ID="}\NormalTok{, }\StringTok{""}\NormalTok{, bound\_genes}\SpecialCharTok{$}\NormalTok{V1)}
\end{Highlighting}
\end{Shaded}

\begin{Shaded}
\begin{Highlighting}[]
\CommentTok{\# Get your significant DEG names}
\NormalTok{sig\_gene\_names }\OtherTok{\textless{}{-}} \FunctionTok{rownames}\NormalTok{(sig)}

\CommentTok{\# Find overlap}
\NormalTok{overlap }\OtherTok{\textless{}{-}} \FunctionTok{intersect}\NormalTok{(sig\_gene\_names, bound\_genes}\SpecialCharTok{$}\NormalTok{V1)}
\FunctionTok{cat}\NormalTok{(}\StringTok{"Number of genes bound AND differentially expressed:"}\NormalTok{, }
    \FunctionTok{length}\NormalTok{(overlap), }\StringTok{"}\SpecialCharTok{\textbackslash{}n}\StringTok{"}\NormalTok{)}
\end{Highlighting}
\end{Shaded}

\begin{verbatim}
## Number of genes bound AND differentially expressed: 32
\end{verbatim}

\begin{Shaded}
\begin{Highlighting}[]
\CommentTok{\# Save overlap list}
\FunctionTok{write.csv}\NormalTok{(}\FunctionTok{data.frame}\NormalTok{(}\AttributeTok{gene=}\NormalTok{overlap),}
  \StringTok{"bound\_and\_DE\_genes.csv"}\NormalTok{)}
\end{Highlighting}
\end{Shaded}

\begin{Shaded}
\begin{Highlighting}[]
\CommentTok{\# PCA Plot {-} save as PNG}
\FunctionTok{png}\NormalTok{(}\StringTok{"PCA\_plot.png"}\NormalTok{, }
    \AttributeTok{width=}\DecValTok{800}\NormalTok{, }\AttributeTok{height=}\DecValTok{600}\NormalTok{)}
\NormalTok{vsd }\OtherTok{\textless{}{-}} \FunctionTok{vst}\NormalTok{(dds, }\AttributeTok{blind=}\ConstantTok{FALSE}\NormalTok{)}
\FunctionTok{plotPCA}\NormalTok{(vsd, }\AttributeTok{intgroup=}\StringTok{"condition"}\NormalTok{)}
\end{Highlighting}
\end{Shaded}

\begin{verbatim}
## using ntop=500 top features by variance
\end{verbatim}

\begin{Shaded}
\begin{Highlighting}[]
\FunctionTok{dev.off}\NormalTok{()}
\end{Highlighting}
\end{Shaded}

\begin{verbatim}
## pdf 
##   2
\end{verbatim}

\begin{Shaded}
\begin{Highlighting}[]
\CommentTok{\# MA Plot {-} save as PNG}
\FunctionTok{png}\NormalTok{(}\StringTok{"MA\_plot.png"}\NormalTok{,}
    \AttributeTok{width=}\DecValTok{800}\NormalTok{, }\AttributeTok{height=}\DecValTok{600}\NormalTok{)}
\FunctionTok{plotMA}\NormalTok{(res, }\AttributeTok{ylim=}\FunctionTok{c}\NormalTok{(}\SpecialCharTok{{-}}\DecValTok{5}\NormalTok{,}\DecValTok{5}\NormalTok{),}
       \AttributeTok{main=}\StringTok{"MA Plot {-} Mutant vs Normal"}\NormalTok{)}
\FunctionTok{dev.off}\NormalTok{()}
\end{Highlighting}
\end{Shaded}

\begin{verbatim}
## pdf 
##   2
\end{verbatim}

\end{document}
